表\ref{tab:examples}按文档ID顺序展示固定测试子集的前四条。初始方案输入正文，最终方案输入导语加正文。

\begin{table}[H]
\centering\small
\caption{标题生成样例}\label{tab:examples}
\begin{tabularx}{\linewidth}{Y Y Y}
\toprule 参考标题 & 初始方案 & 最终方案 \\ \midrule
动物特工: 俄罗斯白鲸和其他那些飞禽走兽 & “间谍”的间谍:\ReportAnimalPronoun{}们为什么会“间谍” & 间谍:\ReportAnimalPronoun{}们不服的“间谍” \\
香港学生罢课 台湾学生声援 朝野关注 & 台湾学生声援香港占中 & 台湾太阳花:陈为廷称香港罢课是“民主” \\
媒体报道揭示上海自由贸易区管理安排 & 北京将上海自贸区列入“国际金融中心” & 上海将设立自贸区 \\
北京两会期间周永康再次受到媒体关注 & 周滨为周永康之子周滨 & 周滨为周永康之子 \\
\bottomrule
\end{tabularx}
\end{table}

样例中有些重复减少，上海自贸区新闻的输出也更贴近事件主题，但动物特工样例仍存在语义不清的问题。为进一步观察数字和事件关系，从按文档ID排序的前8条验证样例中取以下两例。

\begin{quote}\small

\textbf{输入片段}\quad 有报道说,巴黎地区警方前夜扣押12名涉嫌支持上周当地恐怖袭击的嫌疑犯。\par

\textbf{生成标题}\quad 巴黎袭击案:警方前夜逮捕7名嫌疑人\par

\end{quote}

导语明确给出12名嫌疑犯，输出仍写成7名，具体人数被错误生成。

\begin{quote}\small

\textbf{输入片段}\quad 海协会会长陈云林在完成六次陈江会后周三(12月22日)中午离开台湾,他此行在台湾如同过往到访一样,受到不断的抗议与甚高规格的维安保护。\par

\textbf{生成标题}\quad 台湾来鸿:六四后两岸和平大趋势\par

\end{quote}

导语描述完成六次陈江会后离台，输出却出现“六四后”，改变了事件含义。

初始贪心输出中，四字符片段至少重复三次的提示有17条，最终方案降至2条。最终输出中还有18条数字片段未直接出现在完整输入中。这些规则用于定位待复核样本，事实判断还需检查上下文与表达转换。上述案例表明，输入补全和重复约束之后，数字与事件关系仍可能出错。
